
\subsection{0) Unknowns and maps}\label{unknowns-and-maps}

At each integration point \(x_i\) (weight \(w_i\)):

\begin{itemize}
\tightlist
\item
  Basic force vector: \(\bm{q}\in\mathbb{R}^{n_b}\)
\item
  Section deformation vector: \(\mathbf{e}_i \in \mathbb{R}^{n_s}\)
\item
  Force interpolation:
  \(\mathbf{B}_i := \mathbf{B}(x_i)\in\mathbb{R}^{n_s\times n_b}\)
\item
  Interpolated (target) section resultants: \[
  \mathbf{s}^h_i(\bm{q}) = \mathbf{B}_i\,\bm{q} + \mathbf{s}^{p}_i
  \]
\item
  Constitutive (resisting) section resultants: \[
  \mathbf{s}_i(\mathbf{e}_i)
  \]
\item
  Section tangent stiffness: \[
  \mathbf{K}_{s} := \frac{\partial \mathbf{s}_i}{\partial \mathbf{e}_i}\Big|_{\text{current iterate}}
  \]
\item
  Section tangent flexibility: \[
  \mathbf{F}_{s} := \mathbf{K}_{s}^{-1}
  \]
\end{itemize}

Also define the basic deformation vector (given from the global solve)
\(\bm{v}\in\mathbb{R}^{n_b}\).

\subsection{1) Strong residual statement (the coupled nonlinear
equations)}\label{strong-residual-statement-the-coupled-nonlinear-equations}

We solve for \(\bm{q}\) and \(\{\mathbf{e}_i\}\) such that:

\subsubsection{(A) Compatibility (exact kinematic
constraint)}\label{a-compatibility-exact-kinematic-constraint}

\[
\mathbf{r}_v(\{\mathbf{e}_i\}) \;:=\; \bm{v} \;-\; \sum_{i=1}^{n_{IP}} w_i\,\mathbf{B}_i^{T}\mathbf{e}_i \;=\;\mathbf{0}
\tag{1}
\]

\subsubsection{(B) Section force consistency at every integration
point}

I'll write it in the ``FF-sign'' form (target minus resisting):

\[
\mathbf{r}_{s,i}(\bm{q},\mathbf{e}_i)\;:=\;\mathbf{s}^h_i(\bm{q}) - \mathbf{s}_i(\mathbf{e}_i)\;=\;\mathbf{0}
\tag{2}
\]

That's the \emph{strong} residual system: (1) + (2) for all \(i\).

\subsection{2) Newton linearization of the residual
system}\label{newton-linearization-of-the-residual-system}

Let the current iterate be \((\bm{q},\{\mathbf{e}_i\})\). We seek
corrections \((\delta\bm{q},\{\delta\mathbf{e}_i\})\) such that the
linearized residual vanishes.

\subsubsection{2.1 Linearize compatibility residual
(1)}\label{linearize-compatibility-residual-1}

\[
\mathbf{r}_v(\{\mathbf{e}_i+\delta\mathbf{e}_i\})
\approx
\mathbf{r}_v(\{\mathbf{e}_i\})
-
\sum_i w_i\,\mathbf{B}_i^{T}\delta\mathbf{e}_i
\]

Newton condition: \(\mathbf{r}_v + \delta\mathbf{r}_v = \mathbf{0}\),
so:

\[
\sum_i w_i\,\mathbf{B}_i^{T}\delta\mathbf{e}_i
=
\mathbf{r}_v(\{\mathbf{e}_i\})
=
\bm{v}-\sum_i w_i\,\mathbf{B}_i^T\mathbf{e}_i
\tag{3}
\]

\subsubsection{2.2 Linearize section residual (2) at each
IP}\label{linearize-section-residual-2-at-each-ip}

First compute Jacobians:

\begin{itemize}
\tightlist
\item
  \(\dfrac{\partial \mathbf{s}^h_i}{\partial \bm{q}}=\mathbf{B}_i\)
\item
  \(\dfrac{\partial \mathbf{s}_i}{\partial \mathbf{e}_i}=\mathbf{K}_{s}\)
\end{itemize}

So linearization:

\[
\mathbf{r}_{s,i}(\bm{q}+\delta\bm{q},\mathbf{e}_i+\delta\mathbf{e}_i)
\approx
\mathbf{r}_{s,i}(\bm{q},\mathbf{e}_i)
+
\mathbf{B}_i\,\delta\bm{q}
-
\mathbf{K}_{s}\,\delta\mathbf{e}_i
\]

Newton condition \(\mathbf{r}_{s,i}+\delta\mathbf{r}_{s,i}=\mathbf{0}\)
gives:

\[
\mathbf{B}_i\,\delta\bm{q} - \mathbf{K}_{s}\,\delta\mathbf{e}_i
=
-\mathbf{r}_{s,i}(\bm{q},\mathbf{e}_i)
=
-\big(\mathbf{s}^h_i(\bm{q})-\mathbf{s}_i(\mathbf{e}_i)\big)
\tag{4}
\]

Solve (4) for \(\delta\mathbf{e}_i\):

\[
\mathbf{K}_{s}\,\delta\mathbf{e}_i
=
\mathbf{B}_i\,\delta\bm{q}
+ \mathbf{r}_{s,i}
% +
% \big(\mathbf{s}^h_i(\bm{q})-\mathbf{s}_i(\mathbf{e}_i)\big)
\]

\[
\delta\mathbf{e}_i
=
\mathbf{F}_{s}\,\mathbf{B}_i\,\delta\bm{q}
+
\mathbf{F}_{s}\,\big(\mathbf{s}^h_i(\bm{q})-\mathbf{s}_i(\mathbf{e}_i)\big)
= \dd \bm{e}_q + \dd \bm{e}_s
\tag{5}
\]

Now we can answer your ``what is what?'' questions precisely:

\begin{itemize}
\item
  \textbf{What is \(\Delta \mathbf{s}\)?}\\
  In (5), the term driven by \(\delta\bm{q}\) is: \[
  \Delta\mathbf{s}_i := \mathbf{B}_i\,\delta\bm{q}
  \] i.e., the \emph{increment in interpolated/target section forces
  caused by the basic-force correction}.
\item
  \textbf{What is $\mathbf{r}(x)$ (or \(\mathbf{r}_i\)) exactly?}\\
  It is the \emph{Newton deformation correction induced by the current
  section force residual}, i.e. \[
  \boxed{
  \mathbf{r}_i
  :=
  \mathbf{F}_{s}\,\big(\mathbf{s}^h_i(\bm{q})-\mathbf{s}_i(\mathbf{e}_i)\big)
  }
  \tag{6}
  \] No mysticism: it is literally the solution of the \emph{linearized}
  section equation (4) when \(\delta\bm{q}=0\).\\
  Put differently: it's the deformation increment that would (to first
  order) remove the force mismatch at the section.
\end{itemize}


\subsection{\texorpdfstring{3) Eliminate \(\delta\mathbf{e}_i\) to get
the FF force correction
\(\delta\bm{q}\)}{3) Eliminate \textbackslash delta\textbackslash mathbf\{e\}\_i to get the FF force correction \textbackslash delta\textbackslash mathbf\{q\}}}\label{eliminate-deltamathbfe_i-to-get-the-ff-force-correction-deltamathbfq}

Substitute (5) into the linearized compatibility equation (3):

Left-hand side of (3): \[
\sum_i w_i\,\mathbf{B}_i^{T}\delta\mathbf{e}_i
=
\sum_i w_i\,\mathbf{B}_i^{T}
\Big(
\mathbf{F}_{s}\mathbf{B}_i\delta\bm{q}
+
\mathbf{F}_{s}(\mathbf{s}^h_i-\mathbf{s}_i)
\Big)
\]

Group terms:

\[
\Big(\sum_i w_i\,\mathbf{B}_i^{T}\mathbf{F}_{s}\mathbf{B}_i\Big)\delta\bm{q}
+
\sum_i w_i\,\mathbf{B}_i^{T}\mathbf{F}_{s}(\mathbf{s}^h_i-\mathbf{s}_i)
=
\mathbf{r}_v
\tag{7}
\]


Define the \textbf{integrated residual deformation vector} \[
\boxed{
\bm{v}_r := \sum_i w_i\,\mathbf{B}_i^{T}\mathbf{F}_{s}(\mathbf{s}^h_i-\mathbf{s}_i)
}
\tag{9}
\]

Then (7) becomes the Newton system reduced to \(\delta\bm{q}\):

\[
\boxed{
\mathbf{F}\,\delta\bm{q} + \bm{v}_r = \mathbf{r}_v
}
\tag{10}
\]

So the fully general Newton correction is: \[
\boxed{
\delta\bm{q} = \mathbf{F}^{-1}\Big(\mathbf{r}_v - \bm{v}_r\Big)
}
\tag{11}
\]

\subsection{4) The ``pure FF'' update you asked
for}\label{the-pure-ff-update-you-asked-for}

You specifically asked to justify the FF inner-loop update of the form

\[
\delta\bm{q} = \mathbf{F}^{-1}\int \mathbf{B}^T\mathbf{f}_s(\mathbf{s}^h-\mathbf{s}(\mathbf{e}))\,dx
\]

From (11), that corresponds to the case where the right-hand side
\(\mathbf{r}_v\) is taken as zero during the \textbf{inner
force-correction loop} (because the imposed \(\bm{v}\) has already
been enforced/frozen at that stage).

Setting \(\mathbf{r}_v=\mathbf{0}\) in (10):

\[
\mathbf{F}\,\delta\bm{q} + \bm{v}_r = \mathbf{0}
\quad\Rightarrow\quad
\boxed{
\delta\bm{q} = -\mathbf{F}^{-1}\bm{v}_r
}
\tag{12}
\]

Expanding \(\bm{v}_r\) from (9):

\[
\boxed{
\delta\bm{q}
=
-\Big(\int_0^L \mathbf{B}^T\mathbf{F}_s\mathbf{B}\,\dd x\Big)^{-1}
\Big(\int_0^L \mathbf{B}^T\mathbf{F}_s(\mathbf{s}^h-\mathbf{s}(\mathbf{e}))\,\dd x\Big)
}
\tag{13}
\]

That is exactly the FF update (up to sign convention).




\subsection{Where does the ``pure FF'' update come from?}

Your requested FF-looking update is essentially

\[\dd \bm{q} = -\mathbf{F}^{-1}\int_0^L \ForceInterp^{\top} \mathbf{F}_{s}\big(\inters-\bm{s}(\bm{e})\big)\,\dd\reflenx,\]

i.e.

\[F\,\dd \bm{q} = -v_r.\]

That corresponds to the special case of the condensed Newton step
(N-condensed) with

\[\residv[\bm{e}]=0\]

at that stage of the iteration.
